\documentclass[12pt]{article}
\usepackage[T1]{fontenc}
\usepackage[utf8]{inputenc}
\usepackage{lmodern}
\usepackage{textcomp}
\usepackage{enumitem}
\usepackage{geometry}
\geometry{margin=1in}
\begin{document}
\begin{center}
Answer Key - Philosophy - Undergraduate Level Assessment 7 \
\small (Answers are ordered exactly as in the question paper.)
\end{center}

\section*{Multiple Choice}
\begin{enumerate}[label=\arabic*.]
\item \textbf{Answer:} B
\\ \textit{Options:} A) Fallacy of Accident; B) Questionable Analogy; C) Questionable Cause; D) Fallacy of Composition
\\ \textit{Note:} The analogy between the universe and a watch is questionable because it oversimplifies the complexity of the universe and assumes that its existence necessitates a creator in the same way that a watch, which is a manufactured object, requires a watchmaker.
\item \textbf{Answer:} B
\\ \textit{Options:} A) Division; B) False sign; C) False cause; D) Hasty Generalization
\\ \textit{Note:} The term "\_Ad novitatem\_" refers to the fallacy of appealing to novelty, which is classified as a false sign because it asserts that something is better or more valid simply because it is new, rather than based on sound reasoning or evidence.
\item \textbf{Answer:} D
\\ \textit{Options:} A) Equivocation; B) Affirming the Consequent; C) Gambler's Fallacy; D) Inconsistency
\\ \textit{Note:} The phrase "held on a space-available basis" creates an inconsistency with the assurance of confirming a reservation, as it implies that the reservation is not guaranteed, undermining the expectation of a confirmed booking.
\item \textbf{Answer:} C
\\ \textit{Options:} A) False Dilemma; B) Appeal to Pity; C) Begging the question; D) Appeal to Authority; E) Bandwagon Fallacy
\\ \textit{Note:} The statement is an example of begging the question because it assumes the conclusion that men are better drivers without providing independent evidence, instead restating the premise in a different form.
\item \textbf{Answer:} A
\\ \textit{Options:} A) Red herring; B) Slippery slope; C) Either-or fallacy; D) Straw man argument; E) Hasty generalization
\\ \textit{Note:} The correct answer is misleading because the fallacy of accident refers to the erroneous application of a general rule to a specific case, while a red herring is a distraction that diverts attention from the argument at hand; thus, they are distinct logical fallacies and should not be conflated.
\end{enumerate}

\section*{True / False}
\begin{enumerate}[label=\arabic*.]
\setcounter{enumi}{5}
\item \textbf{Answer:} True
\\ \textit{Options:} A) True; B) False
\\ \textit{Note:} Logical behaviorism asserts that mental states can be understood in terms of observable behaviors and the conditions that elicit these behaviors, thus defining mental states as dispositions to act in particular ways.
\item \textbf{Answer:} False
\\ \textit{Options:} A) True; B) False
\\ \textit{Note:} Macedo argues that humanitarian duties are essential for societies in their relations with nonmembers, emphasizing the moral obligation to assist those outside their community.
\item \textbf{Answer:} True
\\ \textit{Options:} A) True; B) False
\\ \textit{Note:} A faulty analogy occurs when the differences between the two things being compared outweigh their similarities, making the conclusion drawn from the comparison invalid.
\item \textbf{Answer:} True
\\ \textit{Options:} A) True; B) False
\\ \textit{Note:} According to Hobbes, the fundamental principle of self-preservation in his social contract theory asserts that individuals have a natural right to defend themselves against threats to their life, especially in a state of war where the absence of a governing authority leads to a condition of constant danger.
\item \textbf{Answer:} True
\\ \textit{Options:} A) True; B) False
\\ \textit{Note:} Anscombe criticizes Sidgwick for endorsing consequentialism because she argues that it fails to adequately account for the moral significance of intentions and the inherent nature of actions, which she believes are essential to ethical evaluation.
\end{enumerate}

\section*{Long Form Response}
\begin{enumerate}[label=\arabic*.]
\setcounter{enumi}{10}
\item \textbf{Answer:} The portrayal of Marduk in "Ludul Bel Nemequi" reflects themes of indifference and detachment through his passive role in the face of human suffering and chaos. Rather than intervening directly to alleviate the protagonist's plight, Marduk embodies a distant deity whose actions seem disconnected from the emotional turmoil experienced by mortals. This representation invites a critical examination of ancient philosophies that grapple with the nature of divine justice and the relationship between gods and humans, suggesting a worldview in which divine beings may not necessarily prioritize human welfare. The implications of this detachment challenge the notion of a benevolent deity and provoke deeper considerations about the existential conditions of human life in the context of divine indifference. Ultimately, this portrayal encourages reflection on the complexities of faith and the human quest for meaning in an indifferent universe.
\\ \textit{Note:} correct because it identifies Marduk's passive and detached portrayal as indicative of a broader philosophical inquiry into the nature of divine justice and the emotional disconnect between deities and human suffering in ancient thought.
\item \textbf{Answer:} To construct a complete truth table for the argument \textasciitilde{}G ⊃ H, \textasciitilde{}H / G, we first note the components: \textasciitilde{}G (not G), H, and \textasciitilde{}H (not H). The truth table reveals that the premises are true when G is false and H is true, while G can still be false under these conditions. This leads to the conclusion that the argument is invalid, as the conclusion G does not logically follow from the premises. A counterexample occurs when G is false and H is true; in this case, the premises \textasciitilde{}G ⊃ H and \textasciitilde{}H are satisfied, but G is false, demonstrating that the argument does not hold in all scenarios. Therefore, the argument is invalid due to this inconsistency.
\\ \textit{Note:} correct because it accurately constructs a truth table that demonstrates the premises can be true while the conclusion is false, thus showing the argument is invalid and providing a specific counterexample where the premises hold but the conclusion does not.
\end{enumerate}

\vfill
\noindent\textit{End of Answer Key}
\end{document}